\begin{blocksection}
\question What would Python display?

\begin{lstlisting}
class Foo(object):
    x = 'bam'

    def __init__(self, x):
        self.x = x

    def baz(self):
        return type(self).x + self.x

class Bar(Foo):
    x = 'boom'

    def __init__(self, x):
        Foo.__init__(self, 'er' + x)

foo = Foo('boo')
\end{lstlisting}
\end{blocksection}

\begin{parts}
\begin{blocksection}
\part \lstinline$>>> Foo.x$

\begin{solution}[2em]
\lstinline$'bam'$
\end{solution}
\end{blocksection}

\begin{blocksection}
\part \lstinline$>>> foo.x$

\begin{solution}[2em]
\lstinline$'boo'$
\end{solution}
\end{blocksection}

\begin{blocksection}
\part \lstinline$>>> foo.baz()$

\begin{solution}[2em]
\lstinline$'bamboo'$
\end{solution}
\end{blocksection}

\begin{blocksection}
\part \lstinline$>>> Foo.baz()$

\begin{solution}[2em]
\lstinline$Error$
\end{solution}
\end{blocksection}

\begin{blocksection}
\part \lstinline$>>> Foo.baz(foo)$

\begin{solution}[2em]
\lstinline$'bamboo'$
\end{solution}
\end{blocksection}

\begin{blocksection}
\part \begin{lstlisting}
>>> bar = Bar('ang')
>>> Bar.x
\end{lstlisting}

\begin{solution}[2em]
\lstinline$'boom'$
\end{solution}
\end{blocksection}

\begin{blocksection}
\part \lstinline$>>> bar.x$

\begin{solution}[2em]
\lstinline$'erang'$
\end{solution}
\end{blocksection}

\begin{blocksection}
\part \lstinline$>>> bar.baz()$

\begin{solution}[2em]
\lstinline$'boomerang'$
\end{solution}
\end{blocksection}
\end{parts}

\begin{questionmeta}
    Suggested Time: 8 min; Difficulty: Medium; Optional
    This existing problem includes class attributes, instance attributes, bound methods, and inheritance. Use the \lstinline{Foo} parts for initial object practice; use the \lstinline{Bar} parts only after inheritance has been introduced. Draw separate locations for class and instance attributes, and make the implicit \lstinline{self} argument explicit when tracing calls.
\end{questionmeta}
